%%% mode: latex
%%% TeX-master: t
%%% End:

\chapter{实验背景}
\label{cha:intro}

深度神经网络(DNN)是一类由多层神经元构成的非线性模型，再函数逼近、模式识别、计算机视觉和自然语言处理等领域具有广泛应用，当其网络规模足够大，前馈神经网络可以逼近任意连续函数。

在实际应用中，神经网络的拟合能力不仅与模型结构有关，还与训练数据规模、网络深度、激活函数等因素密切相关。因此研究这些因素对神经网络拟合能力的影响具有重要意义。

本实验通过构建深度神经网络，对一维非线性震荡特征的函数进行拟合，并研究以下因素对模型性能的影响：
\begin{enumerate}[(1)]
    \item 训练样本数量
    \item 网络结构（深度与宽度）
    \item 激活函数类型  
\end{enumerate}

\section{问题描述}
本实验研究目标函数为
\begin{equation}\label{target_fun}
    y = sin(5\pi x)\log(x+2)
\end{equation}

\section{实验目的}
本实验的主要目标包括：
\begin{enumerate}[(1)]
    \item 利用Pytorch构建前馈神经网络实现一元非线性函数拟合；
    \item 研究训练样本数量对模型预测精度的影响；
    \item 研究网络结构（层数与宽度）对模型性能的影响；
    \item 研究不同激活函数对训练模型和预测精度的影响；
    \item 通过Loss曲线与预测曲线分析模型的欠拟合与过拟合现象；
    \item 掌握深度学习实验设计方法和结果分析方法.
\end{enumerate}